% SUPERSEDED 2026-11-01 by paper/intro_2026-11-01.tex, Dylan's completed first
% draft. Kept for its decision record only; the prose there is the version of
% record and the tone to match.
% Introduction — draft 2026-10-31
% Built from Dylan's own rough draft of the same day, which set the tone and flow.
% Structure follows his four points: (1) why a pangenome over a reference sequence,
% (2) what a pangenome graph is, (3) the privacy need and the attack families,
% (4) PathGuard.
%
% EVERY \cite KEY BELOW IS A PLACEHOLDER. There is no .bib file in this repo and
% not one key resolves. paper/references/REFERENCES.md tracks each one by slot.
%
% OPEN DECISIONS, all Dylan's:
%   A. The dial. This draft states the theorem in tau and names the accuracy
%      ceiling as the reported quantity. The 2026-10-24 direction block says the
%      knob is not tau; at the 0.80 operating point tau is 1 - 10^-2873. The
%      SETTLED 2026-10-31 (Dylan): "keep as drafted". The theorem stays stated
%      in tau and the accuracy ceiling is named as the reported quantity.
%   B. Participation and consent. Left out of 1.3 — the sentence drafted for it had
%      no source we hold. Slot 8.
%   C. 1.4 is longer than PanMixer's closing paragraph. Natural split is before
%      "We then show".
%   D. His own draft merged points 1 and 2 into one paragraph. Kept separate here,
%      because 1.2's material is what 1.3 argues from.

\section{Introduction}

% --- 1.1 Why a pangenome rather than a single reference sequence ---
Linear reference genomes collapse population-level genetic diversity into a single,
arbitrary mosaic sequence, forcing all genetic variation to be characterized as
errors, mismatches, or deviations from that baseline \cite{schneider2017}. The
omission is structural rather than incidental: sequence absent from the backbone has
no coordinates on it, so reads carrying that sequence mismap or fail to map, and the
alleles lost this way are those common in the ancestries the backbone represents
least \cite{ballouz2019}. Pangenome graphs instead represent the genomic landscape of
a population by encoding shared and divergent sequences directly in the topology,
supporting applications including read alignment, genotyping, and variant calling
\cite{eizenga2020,siren2021,ebler2022}. Because non-reference sequence is carried as
an alternative traversal rather than as a deviation from a backbone, alignment bias
is reduced \cite{siren2021} and structural variation is modeled natively rather than
as a symbolic record against that backbone \cite{hickey2024}.

% --- 1.2 What a pangenome graph is, and how an individual sits in it ---
In such a graph, sequence shared across genomes is carried in the nodes and the edges
record which pieces follow which, so each genome embedded in the graph is a path
through them \cite{eizenga2020}. A single node can carry a stretch that many genomes
share, so what distinguishes an individual is not new sequence but which walk through
the shared material is theirs. Paths correspond to individual haplotypes, so a
diploid genome is represented as a pair of haplotype paths through the shared
sequence graph, and linkage disequilibrium (LD) and phasing information are preserved
along each one \cite{liao2023}. The first human pangenome reference was released in
this form, with graph representations produced independently by two construction
methods \cite{liao2023,garrison2024,hickey2024}, and the consortium behind it
continues to add individuals \cite{wang2022}. These graphs are released openly, which
is what makes them useful.

% --- 1.3 What a released path discloses, and the attacks it invites ---
Once a graph has been constructed and released, a new genome can be mapped or
projected onto the graph. The resulting path may itself be privacy-sensitive, even
when the graph and its construction cohort are public: the nodes and edges are
shared, and what is individual is the walk taken through them. That walk is
identifying, because the rare and private alleles strung along it are near-unique to
the one person who carries them, so a path functions as an identifier even though
most of it is shared with everyone \cite{erlich2014}. A path released without
adequate protection is open to attacks of several distinct kinds, among them
rarity-weighted linkage against an external panel \cite{blindenbach2026}, membership
inference \cite{venkatesaramani2023}, and imputation-style reconstruction from the
linkage disequilibrium in a reference panel \cite{browning2018}. Each of these is a
test somebody wrote, so a release that defeats every test on the list says nothing
about the test nobody has written yet. An existing method, PanMixer, addresses this
problem for individuals already in the graph, selecting edits under a utility budget
and reporting that the published attacks no longer succeed \cite{blindenbach2026}.
What is needed alongside it is a release whose disclosure is limited in advance: a
statement, fixed before any attack is run, of how well an adversary can tell one
input path from another.

% --- 1.4 PathGuard ---
We present PathGuard, a randomized genome path release mechanism with the theoretical
guarantee that the output distributions induced by any two input paths lie within a
total variation distance of at most a pre-specified privacy parameter
$\tau\in[0,1)$, so that under equal priors no test can tell which of the two produced
a release with accuracy better than $(1+\tau)/2$. Specifically, we consider as the
input a fixed pangenome graph $G$ constructed from a public cohort
$D=\{g_1,\ldots,g_n\}$ together with a target genome and its mapped path $p_g$
through $G$, which for a diploid genome is a pair of haplotype paths the mechanism
treats as one private input, and our objective is to release a randomized sanitized
path while limiting the statistical distinguishability of the target from individuals
represented in $D$. Two target regimes are kept apart throughout: an external target,
whose genome contributed to neither $G$ nor $D$, and a cohort member, whose own path
the graph already carries and which is the regime in which we compare directly
against PanMixer. The bound covers the released path alone, and not the raw target
genome, the graph, or any other statistic published about the same individual. The
mechanism uses a target-independent cohort path model as a baseline distribution,
with both the output support and the baseline built from $G$ and $D$ alone and never
from the target, and then introduces a bounded preference for paths similar to the
target; the released path is drawn from the resulting tilted distribution rather than
chosen from it. The preference is calibrated so that it cannot make output
distributions induced by different input paths differ by more than the prescribed
total variation distance. The guarantee in $\tau$ covers every pair of input paths and
is kept for that reason, but at the tilt magnitudes at which the mechanism measurably
changes a release $\tau$ is numerically indistinguishable from one, so we
additionally report, per target, an identification-accuracy ceiling obtained from an
exact Kullback--Leibler divergence between two release distributions; that conversion
is an inequality, so the true ceiling is lower than the one we report. Graph-level
differentially private mechanisms for pangenome data already exist \cite{privvg2022},
so what is new here is the released object and the limit attached to it: a path
through a graph held fixed, carrying a bound on identification accuracy stated per
target before any attack is run. We then show that the stated ceiling holds against
every distinguisher we could build, including PanMixer's own published evaluators run
unmodified on our releases; that PathGuard and PanMixer can be placed on one axis, on
one graph build, for individuals already in the graph; and what the choice costs in
utility, measured as movement away from an untilted release rather than as a level
retained. The distance between the distinguishability the bound permits and the
distinguishability any attacker we could build achieves is this paper's central
finding, and it widens as the permitted ceiling rises.
